\subsection{Maximum Principles for Parabolic Equations}\label{sec:parabolic-maximum-principles}

In this section, we shall extend the maximum principle to parabolic equations of the form
\begin{equation}\label{eq:parabolic-maximum-principle-operator}
Lu=-\dfrac{\partial u}{\partial t}+a_{ij}(\bm{x},t)\dfrac{\partial^2u}{\partial x_i\partial x_j}+b_i(\bm{x},t)\dfrac{\partial u}{\partial x_i}+c(\bm{x},t)u=f(\bm{x},t).
\end{equation}


We shall consider solutions defined on $\Omega\times(0,T)$, where $\Omega$ is an (open) domain in $\mathbb{R}^n$ (more general regions in $(\bm{x},t)$ space can be considered). We set $D=\Omega\times(0,T]$, $Q=\Omega\times(0,T)$ and $\Sigma=(\partial\Omega\times[0,T])\cup(\Omega\times\{0\})$. We shall assume that the coefficients are continuous on $\overline D$, and that the matrix $a_{ij}$ is strictly positive definite on $\overline D$. Throughout we shall also assume that $u\in C^2(D)\cap C(\overline D)$. The maximum principle for parabolic equations says that (under appropriate sign conditions on $c$ and $f$) the maximum of $u$ must be on $\Sigma$.

\subsubsection{The Weak Maximum Principle}\label{subsec:parabolic-weak-maximum}

The analogue of Theorem~\ref{thm:weak-maximum-principle} for parabolic equations is the following result.

\begin{Theorem}
\label{thm:parabolic-weak-maximum}
Assume that $\Omega$ is bounded and $Lu\geq0$ ($Lu\leq0$). Moreover, let $c(\bm{x},t)=0$ in $D$. Then the maximum (minimum) of $u$ is achieved on $\Sigma$.
\end{Theorem}

\begin{Proof}
From Theorem~\ref{thm:weak-maximum-principle} and Remark~\ref{rem:weak-maximum-semidefinite}, we already know that the maximum is achieved on $\partial D$. Moreover, if $Lu>0$, then the maximum cannot be achieved on $\Omega\times\{T\}$. For let $\bm{x}\in\Omega\times\{T\}$ be such that $u(\bm{x},T)=\max_{\bm{y}\in\Omega\times\{T\}}u(\bm{y},T)$. Then at the point $\bm{x}$ we have
\begin{equation}\label{eq:parabolic-top-spatial-maximum-inequality}
a_{ij}\dfrac{\partial^2u}{\partial x_i\partial x_j}+b_i\dfrac{\partial u}{\partial x_i}\leq0,
\end{equation}
and since $Lu>0$, we conclude $u_t<0$.

For the general case, let $u_\epsilon=u+\epsilon\exp(-t)$. This yields
\begin{equation}\label{eq:parabolic-weak-positive-perturbation}
Lu_\epsilon=Lu+\epsilon\exp(-t)>0,
\end{equation}
hence
\begin{equation}\label{eq:parabolic-weak-boundary-maximum}
\max_{\overline D}u_\epsilon=\max_\Sigma u_\epsilon.
\end{equation}
The theorem follows by letting $\epsilon\to0$.
\end{Proof}

Obviously, the analogues of Corollaries~\ref{cor:weak-maximum-nonpositive-zero-order} and~\ref{cor:elliptic-comparison} (with $D$ in place of $\Omega$ and $\Sigma$ in place of $\partial\Omega$) also hold. We note that by setting $u=v\exp(\gamma t)$, we find
\begin{equation}\label{eq:parabolic-zero-order-exponential-shift}
Lu=\exp(\gamma t)(Lv-\gamma v).
\end{equation}
Hence, if $Lu\geq0$, then $(L-\gamma)v\geq0$ and vice versa. By choosing $\gamma$ large enough, we can always achieve that $c-\gamma\leq0$. Hence the analogue of Corollary~\ref{cor:elliptic-comparison} holds without any sign condition on $c$.
